\section{Discussion and conclusions}\label{sec:discussion}
The invasion threshold separates the conditions for coexistence from the kind of coexistence that follows. For \(R\le1\), every positive trajectory approaches the predator-free state. At equality the decay of the predator is of order \(t^{-1}\). For \(R>1\), both species persist. The analysis of \(F(x)+C/x\) explains how an increase in prey growth at low density can lead to three feasible equilibria. When the outer equilibria attract and cycles are excluded, the initial populations determine the long-run state through their position relative to the saddle separatrix.

The parameters \(\alpha\) and \(\beta\) influence distinct parts of this mechanism. Increasing \(\alpha\) increases the initial slope of per-capita prey growth and can permit the turning points of the equilibrium curve. Increasing \(\beta\), with other dimensional inputs fixed, lowers the predator invasion ratio and changes the constant term of the coexistence cubic. These statements concern the current response \(mxy/(\beta+y)\), whose interpretation follows directly from that denominator. The displayed predator equation has no separate quadratic mortality term. Adding such a term changes the predator nullcline and requires recalculating both the root and stability conditions.

The Hopf cases show that varying either predation or carrying capacity can destabilize a coexistence equilibrium and produce local attracting oscillations along specified parameter paths. For enrichment this agrees qualitatively with one mechanism discussed by \citet{haque11}; the transition depends on the parameter path and the response function. The numerical loss of stability does not itself predict predator extinction or establish the large-amplitude fate of a periodic branch. Conversely, a periodic trajectory obtained by integration does not prove uniqueness. Our one-orbit case uses a separate global criterion with its hypotheses checked at the chosen inputs.

The literature-informed results show how changing \(r\) affects equilibrium levels and transient shapes under a common illustrative setting. A field prediction would require independent estimates of predation, conversion, interference, mortality and observation scales for each species pair. The stochastic results concern small displacements about a coexistence state in the two stated noise models. The global count of cycles outside the certified regions, continuation far from the codimension-two points, and the zero-trace triple root remain open here.

\clearpage
\appendix
\section{Formula sheet and figure code}\label{sec:figurecode}
\begin{table}[!htbp]\centering
\caption{Quantities used to move from a parameter set to the corresponding stability test.}
\label{tab:formulasheet}
\begin{tabularx}{\linewidth}{@{}lY@{}}\toprule
Quantity & Definition and use\\ \midrule
Invasion & \(R=emk/(d\beta)\). The predator grows when rare at \(E_1\) precisely when \(R>1\).\\
Coexistence & \(Q(x)=x[F(x)-m]+C\) with \(C=d\beta/e\). Keep roots \(C/m<x<k\) and set \(y=emx/d-\beta\).\\
Local stability & \(T=xF'(x)-dF(x)/m\) and \(\Det=-(dF(x)/m)Q'(x)\) at a feasible root.\\
Cycle exclusion & \(A=\alpha k\), \(M=m/r\), \(D=d/r\). The first Dulac bound is \(D\ge M\chi(A)\).\\
Hopf threshold & At a simple equilibrium, \(T=0\) and \(\Det>0\). Check the along-branch crossing derivative and the cubic focal coefficient.\\
\bottomrule\end{tabularx}
\end{table}

The paths in Table~\ref{tab:figurepaths} are relative to the root of the accompanying source package. Its data directory contains machine-readable parameter sets and numerical diagnostics. The figure scripts regenerate the plots from the model and the recorded inputs.

\begin{table}[H]\centering
\caption{Python entry points for all figures in this report.}
\label{tab:figurepaths}
\begin{tabularx}{\linewidth}{@{}cYY@{}}\toprule
Figure & Output under \texttt{figures/} & Python entry point\\ \midrule
\ref{fig:invasion} & \path{num_invasion.pdf} & \path{code/numerical_study.py}\\
\ref{fig:basins} & \path{num_basins.pdf} & \path{code/numerical_study.py}\\
\ref{fig:hopf} & \path{num_hopf.pdf} & \path{code/numerical_study.py}\\
\ref{fig:two} & \path{num_two_parameters.pdf} & \path{code/numerical_study.py}\\
\ref{fig:uniqueorbit} & \path{num_cycle.pdf} & \path{code/numerical_study.py}\\
\ref{fig:normalcycles} & \path{analytical_cycles.pdf} & \path{code/analytical_diagram.py}\\
\ref{fig:bautin-numerical} & \path{num_bautin.pdf} & \path{code/numerical_study.py}\\
\ref{fig:lit-rates} & \path{lit_growth_rates.pdf} & \path{code/literature_cases.py}\\
\ref{fig:noise} & \path{stoch_moments.pdf} & \path{code/stochastic_stability.py}\\
\bottomrule\end{tabularx}
\end{table}

\clearpage
\bibliographystyle{elsarticle-harv}
\bibliography{references}
\clearpage
\section{Research Extensions}\label{sec:extensions}
\noindent\textbf{1. How does the feeding law change coexistence?}\par
We kept the coursework's weak Allee prey growth and changed only the feeding law. All five models preserve positive populations and keep them bounded. In the two prey-dependent models (P1 and P2), positive coexistence fixes prey density independently of carrying capacity. The original foraging-arena response (P0) can support two attracting coexistence states. With ratio-dependent feeding (P4), enrichment lowers the Jacobian trace at every positive coexistence state with positive determinant for the coursework growth law. For a positive growth factor $g$ with $g'>0$, the same decrease holds at a given prey density for every admissible ratio-dependent intake if and only if $(1/g)''\geq0$ there.

\medskip
\noindent\textbf{2. How do prey growth and predator competition interact?}\par
We kept the feeding law fixed and paired three prey growth laws with and without direct predator competition. Logistic growth allows at most one positive equilibrium. Weak and strong Allee growth can each support three. When the three growth laws are calibrated to the same coexistence state at zero competition, positive competition orders predator density along their common branches with positive determinant. Predator density is highest for strong Allee growth, intermediate for weak Allee growth, and lowest for logistic growth. This ranks equilibrium levels even when predator density falls in all three models. In 48 disturbance simulations at the study baseline, 45 recovered. All three failures followed prey removal under strong Allee growth.

\medskip
\noindent\textbf{3. What transfers when the same equation is fitted to other systems?}\par
We kept the ecological equation and fitted its coefficients to training conditions in tumor--T-cell and bacteria--phage datasets. The equation had lower mean held-out error than the CARRGO model on the PBT138 cell line, while the infection-stage model had lower four-hour error in 11 of 12 $\lambda$-phage conditions. In 12 held leukemia conditions, fitting targets alone predicted none of the observed early T-cell increases. Fitting both populations predicted all 12, but mean target root-mean-square error rose from 81,634 to 197,580 cells. The measured trade-off concerns predictions for held conditions from this dataset.

\medskip

\noindent The companion \textit{Appendix --- Research Extensions} gives the figures, proofs, Python calculations, and Lean verification results for each study.
\end{document}
